\documentclass[11pt]{article}
\usepackage{amsmath, amssymb, amsthm}
\usepackage{geometry}
\usepackage{hyperref}
\usepackage{booktabs}
\usepackage{enumitem}
\usepackage[dvipsnames]{xcolor}
\usepackage{tcolorbox}

\geometry{margin=1.1in}
\hypersetup{colorlinks=true, linkcolor=blue!60!black, urlcolor=blue!60!black, citecolor=blue!60!black}

\newcommand{\eps}{\varepsilon}
\newcommand{\code}[1]{\texttt{#1}}
\newcommand{\R}{\mathbb{R}}
\newcommand{\E}{\mathbb{E}}
\newcommand{\Prob}{\mathbb{P}}
\newcommand{\tr}{\operatorname{Tr}}
\newcommand{\dv}{\nabla\cdot}
\newcommand{\KL}{\operatorname{KL}}
\newcommand{\Pc}{\mathcal{P}}
\newcommand{\Nc}{\mathcal{N}}
\newcommand{\Kc}{\mathcal{K}}
\newcommand{\hphi}{\hat\varphi}
\newcommand{\ind}{\mathbf{1}}
\setlist[itemize]{leftmargin=1.4em, itemsep=2pt}
\setlist[enumerate]{leftmargin=1.8em, itemsep=2pt}

\theoremstyle{definition}
\newtheorem{definition}{Definition}[section]
\newtheorem{example}[definition]{Example}
\newtheorem{remark}[definition]{Remark}
\theoremstyle{plain}
\newtheorem{lemma}[definition]{Lemma}
\newtheorem{proposition}[definition]{Proposition}

\newtcolorbox{readbox}{colback=blue!4!white, colframe=blue!50!black, title=Reading Gushchin et al., fonttitle=\bfseries, boxrule=0.5pt, arc=1pt, left=4pt, right=4pt, top=3pt, bottom=3pt}
\newtcolorbox{keybox}{colback=orange!6!white, colframe=orange!70!black, boxrule=0.5pt, arc=1pt, left=4pt, right=4pt, top=3pt, bottom=3pt}

\title{A Primer on Entropic Optimal Transport\\ and the Schr\"odinger Bridge}
\author{Optimal Transport / Schr\"odinger Bridge Project\\
Companion to \code{python/shared/mixture\_sb.py} and \code{paper\_route\_sb.py}}
\date{20 October 2026}

\begin{document}
\maketitle

\begin{abstract}
\noindent This note builds, from scratch, the probabilistic vocabulary needed to read
Gushchin et al., \emph{Building the Bridge of Schr\"odinger} (NeurIPS 2023,
arXiv:2306.10161), Sections 2--3, and to see why our closed-form mixture reference is exact.
Each section introduces only what the next one needs: measures and random variables,
couplings and conditional distributions, entropy, the three optimal transport problems
(Monge, Kantorovich, entropic), weak OT, Brownian motion and the heat kernel, the
Schr\"odinger bridge and its two potentials, and finally the ``choose the potential''
construction. Worked examples are one-dimensional Gaussians, where every object is a
number. Boxes titled \emph{Reading Gushchin et al.} say which paragraph of the paper each
piece unlocks.
\end{abstract}

\tableofcontents

\section{Conventions used throughout}
\label{sec:conv}

Three quantities are written differently in every source. Fix them once.
\begin{itemize}
  \item The noise level of the stochastic process is $\sigma^2$: the reference process is
        $dX_t=\sigma\,dW_t$, i.e.\ Brownian motion with variance $\sigma^2 t$ per coordinate.
  \item Our solver's constraint is $\partial_t\rho+\dv m=\eps\Delta\rho$. Its Fokker--Planck
        equation (Section~\ref{sec:bm}) forces
        \[
          \eps=\tfrac12\sigma^2 .
        \]
  \item Gushchin et al.\ call $\sigma^2$ simply $\epsilon$; Bunne et al.\ (2023) call it $\sigma^2$;
        Janati et al.\ (2020) build it into their regularization constant. So \emph{every $\epsilon$
        in Gushchin's paper is $2\eps$ in our code.} Section~\ref{sec:table} has the full table.
\end{itemize}
The spatial domain is $\R^d$ in the theory and $[0,1]^d$ with no-flux walls in the solver;
Section~\ref{sec:ours} says when the difference matters. The transport cost is always
$c(x,y)=\tfrac12|x-y|^2$.

\section{Measures, densities, random variables}
\label{sec:measures}

\begin{definition}[Probability measure, density]
A probability measure $\mu$ on $\R^d$ assigns to every (measurable) set $A$ a mass
$\mu(A)\in[0,1]$ with $\mu(\R^d)=1$. When $\mu(A)=\int_A\mu(x)\,dx$ for a nonnegative
function $\mu(x)$ we say $\mu$ has a density, and we use the same letter for the measure
and the density. $\Pc(\R^d)$ is the set of probability measures; $\Pc_2(\R^d)$ those
with finite second moment $\int|x|^2d\mu<\infty$.
\end{definition}

Everything our grid solver sees is a density ($\rho_0,\rho_1,\rho_t$), but the theory needs
two measures without densities: the Dirac mass $\delta_a$ (all mass at the point $a$),
and couplings concentrated on a curve (Section~\ref{sec:couplings}).

\begin{definition}[Random variable, law, expectation]
A random variable $X$ with values in $\R^d$ has a \emph{law} (or distribution) $\mu$,
written $X\sim\mu$, meaning $\Prob(X\in A)=\mu(A)$. The expectation of a function of $X$ is
\[
  \E[f(X)]=\int f\,d\mu=\int f(x)\,\mu(x)\,dx .
\]
\end{definition}
The notations $\int f\,d\mu$, $\int f(x)\,d\mu(x)$, $\int f(x)\,\mu(dx)$ and $\E_{X\sim\mu}[f(X)]$
all mean the same thing. A random variable is just a convenient name for ``a point drawn
from $\mu$''; nothing is random in the formulas, which are all integrals.

\begin{definition}[Pushforward]
If $X\sim\mu$ and $T:\R^d\to\R^d$, the law of $T(X)$ is the pushforward $T_\#\mu$, defined by
$(T_\#\mu)(A)=\mu(T^{-1}(A))$. For a smooth invertible $T$ and densities, the change of
variables formula reads $(T_\#\mu)(T(x))\,|\det\nabla T(x)|=\mu(x)$.
\end{definition}

\begin{definition}[Gaussians]
$\Nc(m,\Sigma)$ is the Gaussian with mean $m$ and covariance $\Sigma$; its density at $x$ is
written $\Nc(x;m,\Sigma)$. In 1D, $\Nc(x;m,s^2)=(2\pi s^2)^{-1/2}e^{-(x-m)^2/2s^2}$.
\end{definition}

Two Gaussian identities do most of the work later. Both hold in any dimension $d$, with
$A,B$ symmetric positive-definite $d\times d$ covariance matrices. First, the \emph{product}
of two Gaussian densities in the same variable is a Gaussian density times a constant:
\begin{equation}
  \label{eq:gprod}
  \Nc(y;a,A)\,\Nc(y;b,B)=\Nc(a;b,A+B)\;\Nc(y;m,S),\quad
  S=(A^{-1}\!+\!B^{-1})^{-1},\ \ m=S\,(A^{-1}a+B^{-1}b).
\end{equation}
The new covariance adds \emph{precisions} (inverse covariances), and the new mean is the
precision-weighted average of the two means. In 1D, with $A=\alpha$, $B=\beta$ scalars, this
reads $S=\alpha\beta/(\alpha+\beta)$ and $m=(\beta a+\alpha b)/(\alpha+\beta)$. The
constant $\Nc(a;b,A+B)$ is the integral of the left-hand side over $y$, so it is also an
instance of the second identity. Second, the \emph{convolution} of two Gaussians adds
covariances:
\begin{equation}
  \label{eq:gconv}
  \int \Nc(y;x,A)\,\Nc(x;m,B)\,dx=\Nc(y;m,A+B).
\end{equation}
Probabilistically, \eqref{eq:gconv} says: if $X\sim\Nc(m,B)$ and, given $X=x$, $Y\sim\Nc(x,A)$,
then $Y\sim\Nc(m,A+B)$. ``Adding independent Gaussian noise adds variance.''

\begin{definition}[Product measure]
For $\mu,\nu\in\Pc(\R^d)$, the product $\mu\otimes\nu$ is the measure on $\R^d\times\R^d$
with density $\mu(x)\nu(y)$. It is the law of a pair $(X,Y)$ with $X\sim\mu$, $Y\sim\nu$
and $X,Y$ \emph{independent}.
\end{definition}

\begin{readbox}
Gushchin's Section 2 opens with $\Pc_2(\R^D)$ and calls the two marginals $P$ and $Q$. Our
$\rho_0$ is their $P$, our $\rho_1$ is their $Q$. Measures that are ``absolutely continuous''
are the ones with a density.
\end{readbox}

\section{Couplings and conditional distributions}
\label{sec:couplings}

This is the section the rest of the primer stands on.

\begin{definition}[Joint law, marginals, coupling]
Let $(X,Y)$ be a pair of random vectors. Its \emph{joint law} $\pi\in\Pc(\R^d\times\R^d)$
has density $\pi(x,y)$. The \emph{marginals} are the laws of $X$ and $Y$ alone,
\[
  \pi_0(x)=\int\pi(x,y)\,dy,\qquad \pi_1(y)=\int\pi(x,y)\,dx .
\]
$\pi$ is a \emph{coupling} (or transport plan) of $\rho_0$ and $\rho_1$ if $\pi_0=\rho_0$
and $\pi_1=\rho_1$. The set of all couplings is $\Pi(\rho_0,\rho_1)$.
\end{definition}

$\Pi(\rho_0,\rho_1)$ is never empty: the product $\rho_0\otimes\rho_1$ is in it. At the
other extreme, if $T_\#\rho_0=\rho_1$, the pair $(X,T(X))$ with $X\sim\rho_0$ is a coupling
too. It has no density on $\R^d\times\R^d$: all its mass sits on the graph $\{y=T(x)\}$.
Such a coupling is called \emph{deterministic}, and it is what the Monge problem optimizes
over (Section~\ref{sec:ot}). Think of a coupling as a recipe ``take a particle at $x$ and
send it to $y$'', which may split the mass at $x$ among several $y$.

\begin{definition}[Conditional distribution]
For a coupling with density, the \emph{conditional distribution of $Y$ given $X=x$} is the
probability measure on $\R^d$ with density
\begin{equation}
  \label{eq:cond}
  \pi(y\,|\,x)=\frac{\pi(x,y)}{\rho_0(x)} .
\end{equation}
It is written $\pi(\cdot\,|\,x)\in\Pc(\R^d)$: for each fixed $x$, a probability measure in
the variable $y$.
\end{definition}
Equation \eqref{eq:cond} is a definition, not a theorem; it is the slice of the joint
density at $x$, renormalized to unit mass. It is exactly the ``where does the mass at $x$ go''
recipe above. Three consequences, each used later:
\begin{enumerate}
  \item \textbf{Disintegration.} $\pi(x,y)=\rho_0(x)\,\pi(y|x)$. The joint law is the first
        marginal times the conditional. Books write this as $\pi(dx,dy)=\rho_0(dx)\,\pi(dy|x)$,
        or $\pi=\rho_0\otimes\pi(\cdot|x)$, which also covers the case without densities.
  \item \textbf{Second marginal from the first.}
        \begin{equation}
          \label{eq:mixture}
          \rho_1(y)=\int\rho_0(x)\,\pi(y\,|\,x)\,dx .
        \end{equation}
        If you know where each $x$ sends its mass, the target is determined. This one line
        is the entire idea of Gushchin's construction: choose $\rho_0$ and the conditionals,
        read off $\rho_1$.
  \item \textbf{Conditional expectation.} $\E[Y\,|\,X=x]=\int y\,\pi(y|x)\,dy$, the mean
        destination of the mass at $x$. The map $x\mapsto\E[Y|X=x]$ is called the
        barycentric projection of the plan.
\end{enumerate}
A family $x\mapsto\pi(\cdot|x)$ of probability measures is also called a \emph{Markov kernel}
or \emph{transition kernel}; the words mean ``a conditional distribution for every $x$''.

\begin{example}[Bivariate Gaussian]
\label{ex:bigauss}
Let $(X,Y)$ be jointly Gaussian in 1D with means $\mu_0,\mu_1$, variances $a^2,b^2$ and
covariance $c=\E[(X-\mu_0)(Y-\mu_1)]$. Completing the square in the joint density gives
\begin{equation}
  \label{eq:gausscond}
  \pi(\cdot\,|\,x)=\Nc\!\Big(\mu_1+\frac{c}{a^2}(x-\mu_0),\ \ b^2-\frac{c^2}{a^2}\Big).
\end{equation}
The mean destination is affine in $x$; the spread of the destinations is the same for every
$x$. Two limits: $c=ab$ gives zero conditional variance, so $Y=\mu_1+\frac{b}{a}(X-\mu_0)$
exactly (a deterministic coupling, the Monge map between the two Gaussians); $c=0$ gives
$\pi(\cdot|x)=\Nc(\mu_1,b^2)$ for every $x$, the product coupling. The entropic plan between
Gaussians (Section~\ref{sec:ot}) is this family with $c=C_\sigma$ from Bunne et al.,
sliding from $ab$ at $\sigma=0$ to $0$ as $\sigma\to\infty$.
\end{example}

\begin{readbox}
Gushchin writes a plan as $\pi\in\Pi(P,Q)$ and its conditionals as $\pi(\cdot|x)$, exactly
as here. Whenever the paper says ``$\pi$ is a mixture of $\pi(\cdot|x)$ over $x\sim P$'' it
means the disintegration above. Their sampling procedure for $Q$ is \eqref{eq:mixture}
read as an algorithm: draw $x\sim P$, then $y\sim\pi(\cdot|x)$.
\end{readbox}

\section{Entropy and Kullback--Leibler divergence}
\label{sec:entropy}

\begin{definition}[Differential entropy]
For a density $\mu$, $H(\mu)=-\int\mu\log\mu\,dx$. For a Gaussian,
$H(\Nc(m,\Sigma))=\tfrac12\log\big((2\pi e)^d\det\Sigma\big)$; in 1D, $H=\tfrac12\log(2\pi e\,s^2)$.
\end{definition}
Entropy is large for spread-out densities and goes to $-\infty$ as a density concentrates
on a point. This is the mechanism by which an entropy term forbids deterministic couplings.

\begin{definition}[KL divergence]
For probability measures $\mu,\nu$ with densities,
\[
  \KL(\mu\,\|\,\nu)=\int\mu(x)\log\frac{\mu(x)}{\nu(x)}\,dx .
\]
It is $\ge0$ with equality only when $\mu=\nu$ (Jensen's inequality), it is not symmetric,
and it is $+\infty$ when $\mu$ puts mass where $\nu$ has none (then $\mu$ is not
``absolutely continuous'' with respect to $\nu$, written $\mu\not\ll\nu$).
\end{definition}
In 1D with equal means, $\KL(\Nc(0,a^2)\|\Nc(0,b^2))=\tfrac12\big(\tfrac{a^2}{b^2}-1-\log\tfrac{a^2}{b^2}\big)$.
Note $\KL(\mu\|\nu)=-H(\mu)-\int\mu\log\nu$: KL is negative entropy plus a cross term.

\begin{lemma}[KL splits along conditionals]
\label{lem:chain}
Let $\pi$ and $\nu$ be two measures on $\R^d\times\R^d$ with the same first marginal
$\rho_0$. Then
\[
  \KL(\pi\,\|\,\nu)=\int\rho_0(x)\,\KL\big(\pi(\cdot|x)\,\big\|\,\nu(\cdot|x)\big)\,dx .
\]
\end{lemma}
\begin{proof}
Write $\pi(x,y)=\rho_0(x)\pi(y|x)$ and $\nu(x,y)=\rho_0(x)\nu(y|x)$ inside the logarithm;
the $\rho_0$ factors cancel.
\end{proof}

The case $\nu=\rho_0\otimes\rho_1$, the independent coupling, is the one entropic OT uses:
\begin{equation}
  \label{eq:mi}
  \KL(\pi\,\|\,\rho_0\otimes\rho_1)=\int\rho_0(x)\,\KL\big(\pi(\cdot|x)\,\|\,\rho_1\big)\,dx
  = H(\rho_1)-\int\rho_0(x)\,H\big(\pi(\cdot|x)\big)\,dx .
\end{equation}
This is the \emph{mutual information} of $(X,Y)$: zero for independent variables, infinite
for a deterministic coupling (each $\pi(\cdot|x)$ is a Dirac, entropy $-\infty$). For the
bivariate Gaussian of Example~\ref{ex:bigauss} with correlation $r=c/(ab)$ it equals
$-\tfrac12\log(1-r^2)$. The second form in \eqref{eq:mi} is what makes entropic OT a
``weak'' problem in Section~\ref{sec:weak}: for fixed marginals, the only part that depends
on the plan is the average entropy of the conditionals.

\section{Monge, Kantorovich and entropic optimal transport}
\label{sec:ot}

\begin{definition}[Monge]
Given $\rho_0,\rho_1$, minimize $\int\tfrac12|x-T(x)|^2\rho_0(x)\,dx$ over maps $T$ with $T_\#\rho_0=\rho_1$.
\end{definition}

\begin{definition}[Kantorovich]
Minimize over couplings instead of maps:
\[
  \tfrac12 W_2^2(\rho_0,\rho_1)=\min_{\pi\in\Pi(\rho_0,\rho_1)}\int\!\!\int\tfrac12|x-y|^2\,\pi(x,y)\,dx\,dy
  =\min_{\pi}\ \E_{(X,Y)\sim\pi}\tfrac12|X-Y|^2 .
\]
\end{definition}
Every map gives a coupling, so Kantorovich's value is at most Monge's; it is a relaxation
(a ``weaker'' problem in the sense of fewer constraints). Brenier's theorem says that for
$\rho_0$ with a density the Kantorovich optimum is deterministic, $y=\nabla\psi(x)$ with
$\psi$ convex, so the two values agree. For Gaussians, $T(x)=\mu_1+A(x-\mu_0)$ with
$A=\Sigma_0^{-1/2}(\Sigma_0^{1/2}\Sigma_1\Sigma_0^{1/2})^{1/2}\Sigma_0^{-1/2}$ and
$W_2^2=|\mu_0-\mu_1|^2+\tr\Sigma_0+\tr\Sigma_1-2\tr(\Sigma_0^{1/2}\Sigma_1\Sigma_0^{1/2})^{1/2}$,
the Bures--Wasserstein distance (Janati et al., Section 2).

\begin{definition}[Entropic OT]
For $\sigma^2>0$,
\begin{equation}
  \label{eq:eot}
  \mathrm{OT}_{\sigma^2}(\rho_0,\rho_1)=\min_{\pi\in\Pi(\rho_0,\rho_1)}
  \int\!\!\int\tfrac12|x-y|^2\,d\pi+\sigma^2\,\KL(\pi\,\|\,\rho_0\otimes\rho_1).
\end{equation}
\end{definition}
The penalty rules out deterministic plans (infinite KL) and pulls the plan toward
independence as $\sigma^2$ grows. Replacing $\KL(\pi\|\rho_0\otimes\rho_1)$ by $-H(\pi)$,
or by $\KL(\pi\|\text{Lebesgue})$, changes the objective by a constant that depends only
on $\rho_0,\rho_1$, so all three conventions have the same minimizer. The regularization
constant depends on how the cost is scaled: with $\tfrac12|x-y|^2$ it is $\sigma^2$ as in
\eqref{eq:eot}, with $|x-y|^2$ it is $2\sigma^2$ (Janati's convention). Both are chosen so
that the entropic plan is the one of the Schr\"odinger bridge with noise $\sigma^2$
(Section~\ref{sec:sb}).

\begin{proposition}[Form of the entropic plan]
\label{prop:form}
The minimizer of \eqref{eq:eot} is unique and has the form
\begin{equation}
  \label{eq:plan}
  \pi^*(x,y)=\hphi(x)\,\Kc_{\sigma}(x,y)\,\varphi(y),\qquad
  \Kc_\sigma(x,y)=\exp\!\Big(-\frac{|x-y|^2}{2\sigma^2}\Big),
\end{equation}
for two nonnegative functions $\hphi,\varphi$, the \emph{Schr\"odinger potentials}, which are
determined (up to $\hphi\to\lambda\hphi$, $\varphi\to\varphi/\lambda$) by the two marginal
constraints, the \emph{Schr\"odinger system}
\begin{equation}
  \label{eq:ssystem}
  \rho_0(x)=\hphi(x)\int\Kc_\sigma(x,y)\varphi(y)\,dy,\qquad
  \rho_1(y)=\varphi(y)\int\hphi(x)\Kc_\sigma(x,y)\,dx .
\end{equation}
Conversely, any $\pi$ of the form \eqref{eq:plan} that lies in $\Pi(\rho_0,\rho_1)$ is the
minimizer.
\end{proposition}
Where it comes from: the Lagrangian of \eqref{eq:eot} with multipliers $f(x),g(y)$ for the
two marginal constraints is minimized by
$\pi^*=\exp\big((f(x)+g(y)-\tfrac12|x-y|^2)/\sigma^2\big)\rho_0(x)\rho_1(y)$; set
$\hphi=\rho_0e^{f/\sigma^2}$, $\varphi=\rho_1e^{g/\sigma^2}$. Sinkhorn's algorithm is the
alternating solution of \eqref{eq:ssystem} for $\hphi$ and $\varphi$. The converse is the
key to Gushchin: \emph{the exponential form plus the right marginals is a certificate of
optimality}, because the problem is strictly convex and the form is exactly the first-order
condition.

The conditional distributions of the entropic plan follow from \eqref{eq:cond} and \eqref{eq:ssystem}:
\begin{equation}
  \label{eq:condplan}
  \pi^*(y\,|\,x)=\frac{\Kc_\sigma(x,y)\,\varphi(y)}{\int\Kc_\sigma(x,y')\varphi(y')\,dy'} .
\end{equation}
\begin{keybox}
Read \eqref{eq:condplan} as: \emph{a Gaussian of variance $\sigma^2$ centered at $x$, reweighted by
$\varphi$, renormalized.} The mass at $x$ spreads out by diffusion and is then attracted to
where $\varphi$ is large. Only the terminal potential $\varphi$ enters; $\hphi$ is fixed
afterwards by $\rho_0$. This is the sentence to hold on to while reading Section 3 of the paper.
\end{keybox}

\begin{example}[Gaussians]
If $\varphi$ is the exponential of a quadratic, so is $\Kc_\sigma\varphi$, and the
conditional \eqref{eq:condplan} is Gaussian with a mean affine in $x$: the plan is a bivariate
Gaussian, i.e.\ Example~\ref{ex:bigauss}. Solving \eqref{eq:ssystem} for the covariance $c$
in 1D gives the quadratic $c^2+\sigma^2c=a^2b^2$, hence
$c=C_\sigma=\tfrac12\big(\sqrt{4a^2b^2+\sigma^4}-\sigma^2\big)$. The matrices $D_\sigma$ and
$C_\sigma$ of Bunne et al.\ are the $d$-dimensional version of this quadratic formula.
\end{example}

\begin{readbox}
Section 2 of the paper states entropic OT, remarks that it has a unique minimizer, and
introduces the potentials. Their $\epsilon$ is our $\sigma^2$. If a formula there carries
$\exp(-|x-y|^2/2\epsilon)$, that is $\Kc_\sigma$.
\end{readbox}

\section{Weak optimal transport}
\label{sec:weak}

Kantorovich's cost is a sum over pairs: $\int c\,d\pi=\int\rho_0(x)\big[\int c(x,y)\pi(y|x)dy\big]dx$.
The bracket is \emph{linear} in the conditional $\pi(\cdot|x)$. Weak OT (Gozlan, Roberto,
Samson, Tetali 2017; Backhoff-Veraguas, Beiglb\"ock, Pammer 2019) allows a nonlinear
dependence.

\begin{definition}[Weak OT]
Let $C:\R^d\times\Pc(\R^d)\to\R$ be a cost of a point and a probability measure, convex in
its second argument. The weak OT problem is
\begin{equation}
  \label{eq:weak}
  \min_{\pi\in\Pi(\rho_0,\rho_1)}\int\rho_0(x)\,C\big(x,\pi(\cdot|x)\big)\,dx .
\end{equation}
\end{definition}
If $C(x,\mu)=\int c(x,y)\,d\mu(y)$, this is Kantorovich with cost $c$; nothing changes.
The word ``weak'' refers to the cost seeing the conditional law only through integrals
(``weakly''), not to a weakening of constraints. It is a generalization of Kantorovich,
whereas Kantorovich is a relaxation of Monge.

\begin{proposition}[Entropic OT is weak OT]
By \eqref{eq:mi}, problem \eqref{eq:eot} equals, up to the constant $\sigma^2H(\rho_1)$,
problem \eqref{eq:weak} with
\begin{equation}
  \label{eq:weakcost}
  C(x,\mu)=\int\tfrac12|x-y|^2\,d\mu(y)-\sigma^2H(\mu).
\end{equation}
\end{proposition}
The nonlinear part is the entropy of the conditional: the plan pays for sending the mass at
$x$ to a sharply defined destination.

\begin{lemma}[Gibbs variational principle]
\label{lem:gibbs}
For a function $V$ with $\int e^{-V/\sigma^2}<\infty$,
\[
  \min_{\mu\in\Pc(\R^d)}\Big\{\int V\,d\mu-\sigma^2H(\mu)\Big\}
  \quad\text{is attained only at}\quad
  \mu^*(y)=\frac{e^{-V(y)/\sigma^2}}{\int e^{-V/\sigma^2}} .
\]
\end{lemma}
\begin{proof}
$\int V d\mu-\sigma^2H(\mu)=\sigma^2\KL(\mu\|\mu^*)-\sigma^2\log\int e^{-V/\sigma^2}$, and
$\KL\ge0$ with equality only at $\mu=\mu^*$.
\end{proof}

Weak OT has a dual problem with a single potential $f$ on the target space, and its
optimality condition is pointwise in $x$:
\begin{equation}
  \label{eq:weakopt}
  \pi^*(\cdot\,|\,x)=\arg\min_{\mu\in\Pc(\R^d)}\Big\{C(x,\mu)-\int f\,d\mu\Big\}
  \quad\text{for }\rho_0\text{-almost every }x .
\end{equation}
For the entropic cost \eqref{eq:weakcost} this is Lemma~\ref{lem:gibbs} with
$V(y)=\tfrac12|x-y|^2-f(y)$, giving
\[
  \pi^*(y|x)\ \propto\ e^{-|x-y|^2/2\sigma^2}\,e^{f(y)/\sigma^2}=\Kc_\sigma(x,y)\,\varphi(y),\qquad \varphi:=e^{f/\sigma^2},
\]
which is \eqref{eq:condplan} again. The dual potential of weak OT and the Schr\"odinger
potential are the same object, $\varphi=e^{f/\sigma^2}$. The marginal constraint
$\rho_1=\int\rho_0\,\pi^*(\cdot|x)\,dx$ is the equation that determines $f$ when both
marginals are given. Run backwards, it says: \emph{pick $f$ freely, and $\rho_1$ is whatever
\eqref{eq:mixture} produces.}

\begin{readbox}
The weak-OT paragraph of Section 2 is there only to state \eqref{eq:weakopt}: the optimal
conditional plans are the minimizers of a pointwise problem driven by one potential. If the
notation $C(x,\mu)$ with a measure as second argument is unfamiliar, it is \eqref{eq:weakcost}.
\end{readbox}

\section{Brownian motion, SDEs, the heat kernel}
\label{sec:bm}

\begin{definition}[Brownian motion]
$W_t$, $t\ge0$, is a Brownian motion in $\R^d$ if $W_0=0$, its paths are continuous, and
its increments over disjoint intervals are independent with $W_t-W_s\sim\Nc(0,(t-s)I)$.
\end{definition}
The process $X_t=x+\sigma W_t$ (``Brownian motion with volatility $\sigma$'', written
$dX_t=\sigma\,dW_t$) has $X_t\sim\Nc(x,\sigma^2tI)$. Its \emph{transition density}
\begin{equation}
  \label{eq:heat}
  p_t(x,y)=\Nc(y;\,x,\,\sigma^2tI)=(2\pi\sigma^2t)^{-d/2}\exp\!\Big(-\frac{|x-y|^2}{2\sigma^2t}\Big)
\end{equation}
is the probability density of being at $y$ at time $t$ having started at $x$. In $y$ it
solves the heat equation $\partial_tp=\tfrac{\sigma^2}{2}\Delta p$ with $p_0=\delta_x$; the
same equation holds in $x$ by symmetry. The Gaussian convolution rule \eqref{eq:gconv} is
the Chapman--Kolmogorov identity $p_{s+t}=\int p_s(x,z)p_t(z,y)\,dz$.

At $t=1$, $p_1(x,y)$ is $\Kc_\sigma(x,y)$ of \eqref{eq:plan} up to the constant
$(2\pi\sigma^2)^{-d/2}$. \emph{The entropic kernel is the one-step heat kernel.} This is the
whole reason entropic OT and the Schr\"odinger bridge are the same problem.

\begin{definition}[SDE with drift]
$dX_t=u(t,X_t)\,dt+\sigma\,dW_t$ describes a particle pushed by the velocity field $u$ and
jostled by noise. If $X_0\sim\rho_0$, the law $\rho_t$ of $X_t$ solves the Fokker--Planck
equation
\begin{equation}
  \label{eq:fp}
  \partial_t\rho_t+\dv(\rho_tu)=\tfrac{\sigma^2}{2}\Delta\rho_t .
\end{equation}
\end{definition}
With $m=\rho u$ this is our solver's constraint $\partial_t\rho+\dv m=\eps\Delta\rho$ with
$\eps=\sigma^2/2$. For $u=0$ it is the heat equation.

\begin{definition}[Path measure, Brownian bridge]
The law of the whole trajectory $(X_t)_{t\in[0,1]}$ is a probability measure on the space of
continuous paths; for $dX=\sigma dW$ started from $\rho_0$ it is the \emph{Wiener measure}
$\mathbb{W}^{\sigma}$. Conditioning Brownian motion on $X_0=x$ and $X_1=y$ gives the
\emph{Brownian bridge}, a Gaussian process with
\[
  X_t\sim\Nc\big((1-t)x+ty,\ \sigma^2t(1-t)I\big).
\]
\end{definition}
The bridge interpolates linearly in the mean and bulges in the middle; the term
$\sigma^2t(1-t)I$ in Bunne's $\Sigma_t$ formula is this bulge.

\begin{readbox}
Gushchin's reference process is Brownian motion with volatility $\epsilon$, i.e.\ our
$\sigma^2=\epsilon$ (equivalently $dX=\sqrt\epsilon\,dW$). The joint law of $(X_0,X_1)$
under it has density $\rho_0(x)\,p_1(x,y)$, which they write as $P\otimes$ the heat kernel.
\end{readbox}

\section{The Schr\"odinger bridge}
\label{sec:sb}

\begin{definition}[Dynamic Schr\"odinger bridge]
Among all path measures $\mathbb{P}$ whose time-$0$ and time-$1$ marginals are $\rho_0$ and
$\rho_1$, find the one closest to the Wiener measure:
\[
  \min_{\mathbb{P}:\ \mathbb{P}_0=\rho_0,\ \mathbb{P}_1=\rho_1}\ \KL(\mathbb{P}\,\|\,\mathbb{W}^\sigma).
\]
\end{definition}
Interpretation: Brownian particles start distributed as $\rho_0$; we observe that at time
$1$ they are distributed as $\rho_1$, which is not what diffusion alone would give; the
bridge is the most likely way this happened. Two reformulations make it computable.

\paragraph{Control form (what our solver computes).}
If $\mathbb{P}$ is the law of $dX=u\,dt+\sigma\,dW$ with $X_0\sim\rho_0$, Girsanov's theorem gives
\begin{equation}
  \label{eq:girsanov}
  \KL(\mathbb{P}\,\|\,\mathbb{W}^\sigma)=\frac{1}{2\sigma^2}\,\E\int_0^1|u(t,X_t)|^2\,dt
  =\frac{1}{\sigma^2}\int_0^1\!\!\int\frac{|m|^2}{2\rho_t}\,dx\,dt,\qquad m=\rho_tu .
\end{equation}
Minimizing the right-hand side subject to the Fokker--Planck constraint \eqref{eq:fp} with
$\rho_0,\rho_1$ prescribed is the Benamou--Brenier problem with diffusion,
\begin{equation}
  \label{eq:bb}
  \min_{\rho,m}\int_0^1\!\!\int\frac{|m|^2}{2\rho}\,dx\,dt\quad\text{s.t.}\quad
  \partial_t\rho+\dv m=\eps\Delta\rho,\ \ \eps=\tfrac12\sigma^2,
\end{equation}
which is exactly what the LADMM solver discretizes. The optimal $u$ is the \emph{current}
velocity of the bridge, and $m$ is its momentum.

\paragraph{Static reduction (what the closed forms use).}
Any path measure is its joint endpoint law $\pi(x,y)$ times the conditional law of the path
given the endpoints. Applying Lemma~\ref{lem:chain} on path space,
\[
  \KL(\mathbb{P}\|\mathbb{W}^\sigma)=\KL(\pi\,\|\,\rho_0\otimes p_1)+\int\pi(x,y)\,\KL\big(\mathbb{P}(\cdot|x,y)\,\|\,\text{bridge}_{x\to y}\big)\,dx\,dy ,
\]
where $\rho_0\otimes p_1$ denotes the density $\rho_0(x)p_1(x,y)$ of the endpoints under
Wiener measure. The second term is killed by choosing $\mathbb{P}(\cdot|x,y)$ to be the
Brownian bridge, and the first term equals, up to a constant,
$\sigma^{-2}\big[\int\tfrac12|x-y|^2d\pi+\sigma^2\KL(\pi\|\rho_0\otimes\rho_1)\big]$, the
entropic OT objective \eqref{eq:eot}. Hence:
\begin{keybox}
The Schr\"odinger bridge is the entropic OT plan $\pi^*$ glued to Brownian bridges:
draw $(X_0,X_1)\sim\pi^*$, then run a Brownian bridge between them. Its endpoint law is
$\pi^*$, and its potentials $\hphi,\varphi$ are those of Proposition~\ref{prop:form}.
\end{keybox}

\paragraph{Potentials in time.}
Extend the two potentials to all $t\in[0,1]$ by running the heat equation backwards from
$\varphi$ and forwards from $\hphi$:
\begin{equation}
  \label{eq:phit}
  \varphi_t(x)=\int p_{1-t}(x,y)\,\varphi(y)\,dy,\qquad
  \hphi_t(y)=\int\hphi(x)\,p_t(x,y)\,dx ,
\end{equation}
so $\partial_t\varphi_t=-\tfrac{\sigma^2}{2}\Delta\varphi_t$ with $\varphi_1=\varphi$, and
$\partial_t\hphi_t=\tfrac{\sigma^2}{2}\Delta\hphi_t$ with $\hphi_0=\hphi$.

\begin{proposition}[Marginals and drift of the bridge]
\label{prop:bridge}
With $\varphi_t,\hphi_t$ as in \eqref{eq:phit} (and $p_1$ in place of $\Kc_\sigma$ in
\eqref{eq:plan}),
\begin{equation}
  \label{eq:rhot}
  \rho_t=\varphi_t\,\hphi_t,\qquad u(t,x)=\sigma^2\nabla\log\varphi_t(x),\qquad
  m=\rho_tu=\sigma^2\,\hphi_t\nabla\varphi_t .
\end{equation}
\end{proposition}
\begin{proof}[Check]
At $t=0$, $\varphi_0=\int p_1\varphi$ and $\hphi_0=\hphi$, so $\rho_0=\hphi\varphi_0$ is the
first Schr\"odinger equation \eqref{eq:ssystem}; similarly at $t=1$. For $0<t<1$, differentiate
the product: $\partial_t(\varphi_t\hphi_t)=-\tfrac{\sigma^2}{2}\hphi_t\Delta\varphi_t+\tfrac{\sigma^2}{2}\varphi_t\Delta\hphi_t$,
while $\dv(\rho_tu)=\sigma^2\dv(\hphi_t\nabla\varphi_t)$ and
$\tfrac{\sigma^2}{2}\Delta(\varphi_t\hphi_t)=\tfrac{\sigma^2}{2}(\hphi_t\Delta\varphi_t+2\nabla\varphi_t\!\cdot\!\nabla\hphi_t+\varphi_t\Delta\hphi_t)$.
Expanding $\dv(\hphi_t\nabla\varphi_t)=\nabla\hphi_t\!\cdot\!\nabla\varphi_t+\hphi_t\Delta\varphi_t$
shows that \eqref{eq:fp} holds term by term. So $\rho_t$ is a Fokker--Planck solution with
the right endpoints and drift $u$, and $\rho_t$ is also the time-$t$ marginal of the
bridge-glued plan, which is the mixture $\int\pi^*(x,y)\,\Nc(\cdot;(1-t)x+ty,\sigma^2t(1-t))$.
\end{proof}
The drift is a gradient of a log: Schr\"odinger bridges are always gradient flows of
$\log\varphi_t$, in the same way that Brenier maps are gradients of convex functions. In
our convention $m=2\eps\,\rho_t\nabla\log\varphi_t$.

\begin{readbox}
This is the ``dynamic SB'' part of Section 2. The paper's $\varphi^*$ and $\hat\varphi^*$
(or $\phi,\hat\phi$; check the pdf) are our $\varphi,\hphi$; the optimal drift
$\epsilon\nabla\log\varphi_t$ is \eqref{eq:rhot} with $\epsilon=\sigma^2$. Their statement
that the optimal process ``is the reference process with an added drift'' is \eqref{eq:girsanov}.
\end{readbox}

\section{Gushchin's construction: choose the potential}
\label{sec:gushchin}

Normally $\rho_0,\rho_1$ are given and \eqref{eq:ssystem} must be solved numerically for
$\hphi,\varphi$. The construction runs this backwards.

\begin{proposition}[Constructed pairs, Gushchin et al.\ Thm.\ 3.2 in substance]
\label{prop:construct}
Fix $\rho_0$ and any nonnegative $\varphi$ with $\varphi_0:=\int p_1(\cdot,y)\varphi(y)\,dy$
finite. Define
\begin{equation}
  \label{eq:construct}
  \hphi:=\frac{\rho_0}{\varphi_0},\qquad
  \pi^*(y|x):=\frac{p_1(x,y)\varphi(y)}{\varphi_0(x)},\qquad
  \rho_1(y):=\int\rho_0(x)\,\pi^*(y|x)\,dx=\varphi_1(y)\,\hphi_1(y),
\end{equation}
where $\varphi_1=\varphi$ by the definition \eqref{eq:phit}, so that the endpoints read
$\rho_0=\varphi_0\hphi_0$ and $\rho_1=\varphi_1\hphi_1$, the $t=0$ and $t=1$ cases of \eqref{eq:rhot}.
Then $\pi^*(x,y)=\rho_0(x)\pi^*(y|x)=\hphi(x)p_1(x,y)\varphi(y)$ is the entropic plan between
$\rho_0$ and this $\rho_1$, and the Schr\"odinger bridge between them has marginals and
drift given by \eqref{eq:rhot}.
\end{proposition}
\begin{proof}
$\pi^*$ has the exponential form, its first marginal is $\rho_0$ by construction (the
conditionals are normalized), and its second marginal is $\rho_1$ by definition. The
converse part of Proposition~\ref{prop:form} does the rest.
\end{proof}
The price is that $\rho_1$ is an output, not an input. The gain is that the drift
$\sigma^2\nabla\log\varphi_t$ and every $\rho_t$ are exact, with no optimization anywhere.

\paragraph{Gaussian-mixture potentials.}
Take
\begin{equation}
  \label{eq:mixphi}
  \varphi(y)=\sum_{k=1}^K w_k\,\Nc(y;\mu_k,\Sigma_k),\qquad w_k>0 .
\end{equation}
The number of components $K$ and the triples $(w_k,\mu_k,\Sigma_k)$ are \emph{free design
parameters}, chosen by hand; nothing in the problem fixes them, and they are not the modes and
covariances of $\rho_1$ (compare $\mu_\pm=\pm a$ with the actual centres $\bar y_\pm$ in
Example~\ref{ex:1d}). They are the input; the marginals are the output. Gushchin pick theirs
``to make the constructed distributions visually pleasant'' (their Appendix~B); ours are the
$(w,b,Q)$ and $(u,a,P)$ arguments of \code{mixture\_sb.py}.
Gushchin call this an LSE (log-sum-exp) potential because $\log\varphi$ is a log-sum-exp of
quadratics. Everything becomes explicit through the two Gaussian identities of
Section~\ref{sec:measures}:
\begin{itemize}
  \item \textbf{Time potential.} By the convolution rule \eqref{eq:gconv},
        $\varphi_t(x)=\sum_kw_k\,\Nc\big(x;\mu_k,\Sigma_k+\sigma^2(1-t)I\big)$:
        running the heat equation backwards inflates each component's covariance.
  \item \textbf{Drift.} $u=\sigma^2\nabla\log\varphi_t$ is a weighted average of the Gaussian
        gradients,
        \begin{align*}
          u(t,x)&=-\sigma^2\sum_k\omega_k(t,x)\,\big(\Sigma_k+\sigma^2(1-t)I\big)^{-1}(x-\mu_k),\\
          \omega_k(t,x)&\propto w_k\,\Nc\big(x;\mu_k,\Sigma_k+\sigma^2(1-t)I\big),
        \end{align*}
        with the responsibilities $\omega_k$ summing to one. Each component pulls toward its
        centre; the pull is strongest where that component dominates.
  \item \textbf{Conditional plan.} $\pi^*(y|x)\propto\varphi(y)\,\Nc(y;x,\sigma^2I)$ is a
        mixture of products of two Gaussians in $y$, so by the product rule \eqref{eq:gprod}
        with $A=\Sigma_k$, $B=\sigma^2I$, i.e.\
        $S_k=(\Sigma_k^{-1}+\sigma^{-2}I)^{-1}$ and $m_k(x)=S_k(\Sigma_k^{-1}\mu_k+\sigma^{-2}x)$,
        \begin{equation}
          \label{eq:mixcond}
          \pi^*(y|x)=\sum_k\omega_k(0,x)\,\Nc\big(y;m_k(x),S_k\big).
        \end{equation}
        The mass at $x$ is split among the $K$ components with $x$-dependent weights, and
        each part lands in a Gaussian around a point between $x$ and $\mu_k$.
  \item \textbf{Target.} $\rho_1=\int\rho_0(x)\pi^*(\cdot|x)\,dx$ is sampled by
        $x\sim\rho_0$, then $y\sim\pi^*(\cdot|x)$ (Gushchin's Prop.\ 3.3). Its density is a
        $d$-dimensional integral, closed form only when $\rho_0$ is itself Gaussian or a
        Gaussian mixture, which is the case treated next.
\end{itemize}

\paragraph{Dictionary: the same objects in the paper's notation.}
Section~3.3 of the paper looks different from \eqref{eq:mixphi}--\eqref{eq:mixcond} for three
reasons: they parametrize the \emph{dual} potential $f^*$ rather than the Schr\"odinger potential
$\varphi$; they write Gaussians \emph{unnormalized}, through
$Q(y\,|\,b,A):=\exp\!\big(-\tfrac12(y-b)^{\!\top}A\,(y-b)\big)$, a bell of precision $A$ with no
$(2\pi)^{-D/2}\det^{-1/2}$ in front; and they parametrize each component by its precision
rather than its covariance. To keep one index, we write their component index $n$ as our $k$,
their $\epsilon$ as our $\sigma^2$, and their weight $w_n$ as $c_k$, because our $w_k$ is
already taken. Two of their names clash with ours and mean different things: their
$\Sigma_n$ and $\mu_n(x)$ are the covariance and mean of the \emph{conditional} plan, our
$S_k$ and $m_k(x)$, not the component parameters $\Sigma_k,\mu_k$ of \eqref{eq:mixphi}.
\begin{enumerate}
  \item \textbf{Dual vs.\ Schr\"odinger potential (their (14), Cor.\ 3.4).} Their conditional
        plan is $\pi^*(y|x)\propto\exp\big((f^*(y)-\tfrac12|x-y|^2)/\epsilon\big)$. With
        \[
          \varphi(y)=\exp\big(f^*(y)/\epsilon\big)
        \]
        this is exactly $\Kc_\sigma(x,y)\varphi(y)$, i.e.\ our \eqref{eq:condplan}. Corollary~3.4
        says precisely this: $\exp(f^*/\epsilon)$ is the Schr\"odinger potential.
  \item \textbf{The LSE potential (their (15)).}
        $f^*(y)=\epsilon\log\sum_kc_k\,Q(y|b_k,\epsilon^{-1}A_k)$ exponentiates to
        \[
          \varphi(y)=\sum_kc_k\,Q\big(y\,|\,b_k,\epsilon^{-1}A_k\big)
                   =\sum_kc_k\exp\!\Big(-\tfrac12(y-b_k)^{\!\top}\tfrac{A_k}{\epsilon}(y-b_k)\Big),
        \]
        a sum of bells with centres $b_k$ and precisions $A_k/\epsilon$. This is \eqref{eq:mixphi}
        with
        \[
          \mu_k=b_k,\qquad \Sigma_k^{-1}=A_k/\epsilon,\qquad
          w_k=c_k\,(2\pi)^{D/2}\det(\Sigma_k)^{1/2},
        \]
        the last line converting unnormalized bells to normalized densities. Only ratios of the
        weights matter (an overall factor in $\varphi$ is absorbed into $\hphi$).
  \item \textbf{Proposition 3.3 is \eqref{eq:mixcond}.} Substituting $\Sigma_k^{-1}=A_k/\epsilon$
        and $\sigma^{-2}=\epsilon^{-1}$ into our $S_k$, $m_k$:
        \begin{align*}
          S_k&=(\Sigma_k^{-1}+\sigma^{-2}I)^{-1}=\epsilon\,(A_k+I)^{-1}
            &&\text{(their $\Sigma_n$)},\\
          m_k(x)&=S_k(\Sigma_k^{-1}\mu_k+\sigma^{-2}x)=(A_k+I)^{-1}(A_kb_k+x)
            &&\text{(their $\mu_n(x)$)}.
        \end{align*}
        Their mixture weights are $\gamma_k=\tilde c_k/\sum_j\tilde c_j$ with
        \[
          \tilde c_k=c_k\,(2\pi)^{D/2}\det(S_k)^{1/2}\,Q\big(x\,|\,b_k,\ \epsilon^{-1}I-\epsilon^{-2}S_k\big)
          \qquad\text{(their $\tilde w_n$)},
        \]
        and these are our responsibilities $\omega_k(0,x)\propto w_k\Nc(x;\mu_k,\Sigma_k+\sigma^2I)$.
        To see it, note the precision in their $Q$ is
        \[
          \epsilon^{-1}I-\epsilon^{-2}S_k=\epsilon^{-1}(A_k+I)^{-1}A_k
          =(\epsilon A_k^{-1}+\epsilon I)^{-1}=(\Sigma_k+\sigma^2I)^{-1},
        \]
        so $Q(x|b_k,\cdot)$ is the unnormalized version of $\Nc(x;\mu_k,\Sigma_k+\sigma^2I)$; and
        \[
          \det(S_k)^{1/2}=\epsilon^{D/2}\det(A_k+I)^{-1/2}
          \propto\det(\Sigma_k)^{1/2}\det(\Sigma_k+\sigma^2I)^{-1/2}
        \]
        supplies exactly the normalizations that turn $c_k$ into $w_k$ and the bell into a
        density, up to the common factor $(2\pi\epsilon)^{D/2}$ that cancels in $\gamma_k$. Their
        normalizer $Z_x$ is our $\varphi_0(x)$ up to the same constant. Their appendix derivation
        (completing the square in the exponent, their (21) and following) is the product rule
        \eqref{eq:gprod} done by hand.
  \item \textbf{Corollary 3.5 is the drift of \eqref{eq:rhot}.} Write $S_k^t:=\epsilon\big((1-t)A_k+I\big)^{-1}$
        (their $\Sigma_n^t$, with $A_n^t=(1-t)A_n$). Their (16) reads
        \[
          v^*(x,t)=\epsilon\nabla\log\sum_kc_k\det(S_k^t)^{1/2}\,
          Q\Big(x\,\Big|\,b_k,\ \tfrac{1}{\epsilon(1-t)}I-\tfrac{1}{\epsilon^2(1-t)}S_k^t\Big).
        \]
        The precision simplifies the same way,
        \[
          \tfrac{1}{\epsilon(1-t)}\big(I-((1-t)A_k+I)^{-1}\big)=\epsilon^{-1}\big((1-t)A_k+I\big)^{-1}A_k
          =(\Sigma_k+\sigma^2(1-t)I)^{-1},
        \]
        and $\det(S^t_k)^{1/2}$ supplies the normalization, so the sum is $\varphi_t(x)$ of the
        \textbf{Time potential} bullet and $v^*=\sigma^2\nabla\log\varphi_t=u$. (Their (12) is
        \eqref{eq:phit}--\eqref{eq:rhot} with the heat kernel written as $\Nc(y|x,(1-t)\epsilon I)$.)
\end{enumerate}
Two things in Section~3.3 are \emph{not} just notation. First, they only require the eigenvalues of
$A_k$ to lie in $(-1,\infty)$, so a component may have $A_k\preceq0$: a flat or even growing
bell, not a Gaussian. Everything still works because every formula above uses $(A_k+I)^{-1}$,
never $A_k^{-1}$: the heat kernel's precision $\epsilon^{-1}I$ dominates and $S_k\succ0$.
Our Gaussian-mixture $\varphi$ is the $A_k\succ0$ sub-case; to cover the rest, read
$\Sigma_k^{-1}:=A_k/\epsilon$ as the primary object in \eqref{eq:mixcond} and never form
$\Sigma_k$ itself. Second, Proposition~3.3 produces \emph{samples} of $P_1=\pi^*_1$ and nothing
else; the paper says explicitly that $P_1$ is in general not a Gaussian mixture even when $P_0$
is. Section~\ref{sec:ours} gives two recipes that close that gap: Recipe~A follows the paper and adds one convolution; Recipe~B, previewed below, makes everything explicit.

\begin{example}[One-dimensional worked case]
\label{ex:1d}
Let $\rho_0=\Nc(0,s^2)$ and $\varphi(y)=\tfrac12\Nc(y;-a,\tau_0^2)+\tfrac12\Nc(y;a,\tau_0^2)$.
Write $\tau^2=\tau_0^2+\sigma^2$. Then
\[
  \varphi_0(x)=\tfrac12\Nc(x;-a,\tau^2)+\tfrac12\Nc(x;a,\tau^2),\qquad
  \omega_\pm(x)=\frac{\Nc(x;\pm a,\tau^2)}{\Nc(x;-a,\tau^2)+\Nc(x;a,\tau^2)}=\frac{1}{1+e^{\mp2ax/\tau^2}},
\]
so a particle at $x>0$ goes to the right branch with probability $\omega_+(x)$ and to the
left branch with probability $\omega_-(x)=1-\omega_+(x)$ (the ``cross-branch'' mass). Given
the branch, $y\sim\Nc\big(m_\pm(x),S\big)$ with
$S=\big(\tau_0^{-2}+\sigma^{-2}\big)^{-1}$ and $m_\pm(x)=S(\pm a/\tau_0^2+x/\sigma^2)$.
The drift that does the splitting is exact: $u=\sigma^2\partial_x\log\varphi_t$ with
$\varphi_t=\tfrac12\Nc(x;-a,\tau_0^2+\sigma^2(1-t))+\tfrac12\Nc(x;a,\tau_0^2+\sigma^2(1-t))$.
With $a=1$, $\tau_0=0.5$, $\sigma^2=0.1$ the cross-branch probability $\omega_-(x)$ is
about $3\%$ at $x=0.6$ and $0.3\%$ at $x=1$; shrinking $a$ or raising $\sigma^2$ increases it.

The target, however, is \emph{not} closed form here: $\rho_1(y)=\int\rho_0(x)\pi^*(y|x)\,dx$
contains $\omega_\pm(x)$, whose denominator $\varphi_0(x)$ is a sum of two Gaussians and does
not cancel against anything. That integral is what Gushchin sample (Prop.\ 3.3). It is also
what \code{paper\_route\_sb.py} evaluates, not point by point but as $\rho_1=\varphi\,\hphi_1$
with $\hphi_1$ obtained by one Gaussian convolution of $\hphi=\rho_0/\varphi_0$ on a fine grid,
which is cheap and spectrally accurate; the ring example of
\code{candidate\_3d\_problems.tex} is built that way with a $\rho_0$ that is no mixture at all.

\emph{The closed-form variant (ours, not the paper's; Recipe~B of Section~\ref{sec:ours}).}
The paper route already gives a usable grid reference with one numerical step (Recipe~A). The
variant here removes that step and makes $\rho_1$ explicit, at the price of a mixture
$\rho_0$. The obstruction is the denominator
$\varphi_0$ in $\hphi=\rho_0/\varphi_0$: a Gaussian divided by a sum of Gaussians is not a
Gaussian, so $\hphi$, $\hphi_1$ and $\rho_1=\varphi\,\hphi_1$ are not mixtures. The second
Schr\"odinger equation says $\rho_1$ is the \emph{product} $\varphi\,\hphi_1$, and products of
mixtures are mixtures by \eqref{eq:gprod}. So the fix is to change which two objects are the
inputs. Gushchin choose $(\rho_0,\varphi)$ and $\hphi$ is forced. We choose $(\hphi,\varphi)$,
both Gaussian, and \emph{both} marginals are outputs. This is legitimate by the converse in
Proposition~\ref{prop:form}: $\pi=\hphi\,p_1\,\varphi$ is the entropic-optimal plan between
its own two marginals, whatever $\hphi,\varphi$ are, provided it has mass one. Concretely, keep
$\varphi$ as above and take a single Gaussian
\[
  \hphi(x)=Z^{-1}\Nc(x;0,\hat s^2),\qquad Z:=\Nc(a;0,\hat s^2+\tau^2),
\]
where the role of $Z$ is only to make the total mass one (step 2 below). Four steps:
\begin{enumerate}
  \item \textbf{Flow $\varphi$ back to $t=0$.} By \eqref{eq:gconv} each variance grows by
        $\sigma^2$: $\varphi_0(x)=\tfrac12\Nc(x;-a,\tau^2)+\tfrac12\Nc(x;a,\tau^2)$,
        $\tau^2=\tau_0^2+\sigma^2$, as before.
  \item \textbf{Multiply: $\rho_0=\hphi\,\varphi_0$.} Apply \eqref{eq:gprod} term by term
        with $A=\hat s^2$ (mean $0$) and $B=\tau^2$ (mean $\pm a$):
        \[
          \Nc(x;0,\hat s^2)\,\Nc(x;\pm a,\tau^2)
          =\underbrace{\Nc(\pm a;0,\hat s^2+\tau^2)}_{=\,Z\ \text{for both signs}}\;
           \Nc(x;\bar x_\pm,v),\qquad
          \bar x_\pm=\pm\frac{a\,\hat s^2}{\hat s^2+\tau^2},\quad v=\frac{\hat s^2\tau^2}{\hat s^2+\tau^2}.
        \]
        The mean is the precision-weighted average of $0$ and $\pm a$; the pull toward the
        origin is stronger the narrower $\hphi$ is. The two constants are equal by symmetry
        and equal $Z$ by definition, so $Z^{-1}$ cancels them:
        \[
          \rho_0=\tfrac12\Nc(x;\bar x_+,v)+\tfrac12\Nc(x;\bar x_-,v),
        \]
        a probability density.
  \item \textbf{Flow $\hphi$ forward to $t=1$.} $\hphi_1(y)=Z^{-1}\Nc(y;0,\hat s^2+\sigma^2)$.
  \item \textbf{Multiply: $\rho_1=\varphi\,\hphi_1$.} Same rule with variances $\tau_0^2$
        and $\hat s^2+\sigma^2$. The constant that appears is
        $\Nc(a;0,\hat s^2+\sigma^2+\tau_0^2)=Z$, the \emph{same} number, because
        $\tau^2=\tau_0^2+\sigma^2$. This is not a coincidence: both constants are the total
        mass of the plan, counted once at $t=0$ and once at $t=1$. Hence
        \[
          \rho_1=\tfrac12\Nc(y;\bar y_+,v')+\tfrac12\Nc(y;\bar y_-,v'),\qquad
          \bar y_\pm=\pm\frac{a(\hat s^2+\sigma^2)}{\hat s^2+\sigma^2+\tau_0^2},\quad
          v'=\frac{(\hat s^2+\sigma^2)\tau_0^2}{\hat s^2+\sigma^2+\tau_0^2}.
        \]
\end{enumerate}
The picture: the plan $\pi^*$ is two bivariate Gaussians, a $+$ branch and a $-$ branch, each
of mass exactly $\tfrac12$ at $t=0$ and $t=1$ and hence at every $t$ (Section~\ref{sec:ours}).
Each branch is a single Gaussian bridge, so $\rho_t$ is always a sum of two Gaussians whose
centres slide from $\bar x_\pm$ to $\bar y_\pm$ and whose variances go from $v$ to $v'$. At
$t=0$ the narrow $\hphi$ pulls the centres together and $\rho_0$ is one bump; at $t=1$ the
broader $\hphi_1$ pulls less and $\rho_1$ is mildly bimodal. The cross-branch weight
$\omega_\pm(x)$ of the paper route still exists but is never needed: it is already summed
into the two branch masses. Gushchin's Proposition~3.3 is not replaced, it is bypassed; the
paper never needs the density of $\rho_1$, the grid solver does.
With $a=1$, $\hat s=0.3$, $\tau_0=0.5$, $\sigma^2=0.1$:
$\rho_0$ has centres $\pm0.20$ and standard deviation $0.27$, a single broad bump; $\rho_1$
has centres $\pm0.43$ and standard deviation $0.33$, a flat-topped, mildly bimodal density.
\end{example}

\begin{readbox}
Map of the paper's Section 3 (arXiv v2 numbering). 3.1, Theorem 3.1: the construction for a
general weak cost. 3.2, Theorem 3.2 with its eq.\ (14): Proposition~\ref{prop:construct}, the
conditional plan $\propto\exp((f^*-c)/\epsilon)$, i.e.\ \eqref{eq:condplan} with
$\varphi=e^{f^*/\epsilon}$. 3.3, eq.\ (15): the LSE potential, our \eqref{eq:mixphi};
Proposition 3.3: the Gaussian-mixture conditionals, our \eqref{eq:mixcond}, used to sample
$P_1$. 3.4, Corollary 3.4: $e^{f^*/\epsilon}$ is the Schr\"odinger potential; Corollary 3.5,
eq.\ (16): the drift, our $u=\sigma^2\nabla\log\varphi_t$. The dictionary above translates each
symbol. Everything from Section 4 on is about the high-dimensional benchmark pairs and neural
solvers and can be skipped.
\end{readbox}

\section{Two recipes for an exact grid reference}
\label{sec:ours}

\paragraph{Why we need a density at all.}
The paper never evaluates $P_1(y)$ at a point, and never needs to. Its solvers are neural: they
take two bags of samples $\{x_i\}\sim P_0$, $\{y_j\}\sim P_1$ as input and return a sampler
(a network producing draws from $\hat\pi(\cdot|x)$, or a learned drift $\hat v$ integrated as an
SDE); they are scored by sample-based moment metrics (BW$_2^2$-UVP). Proposition~3.3, a fast
exact sampler, is therefore all they need. Our solver is a grid discretization of
\eqref{eq:bb}: its input is $\rho_0,\rho_1$ as arrays of cell values, since the marginal
constraints are imposed pointwise, and its output is $\rho_t,m$ at every cell and time step.
Histogramming samples onto the grid would inject Monte Carlo noise of order
$N_{\mathrm{cell}}^{-1/2}$ per cell, far above the $O(h^2)$ discretization error the convergence
suites are designed to measure; and the error norms of $\rho_t$ and $m$ need the exact
$\rho_t$ at every cell. So we need $\rho_0$, $\rho_1$, $\rho_t$ and $m$ as functions we can
evaluate anywhere, to an accuracy well below $10^{-3}$. Two recipes do this. Neither is in
the paper as such; the first follows it closely, the second departs in one choice.

\paragraph{What both recipes share.}
Both rest on one fact, the converse in Proposition~\ref{prop:form} (the paper's Theorem~3.2):
\begin{keybox}
Take \emph{any} two positive functions $\hphi,\varphi$ on $\R^d$ and set
$\pi(x,y)=\hphi(x)\,\Kc_\sigma(x,y)\,\varphi(y)$, scaled to total mass one. Then $\pi$ is the
entropic-optimal plan, and the Schr\"odinger bridge, between its own two marginals. By
\eqref{eq:phit}--\eqref{eq:rhot} the bridge is
\[
  \rho_t=\varphi_t\,\hphi_t,\qquad u=\sigma^2\nabla\log\varphi_t,\qquad m=\rho_tu,
\]
with $\varphi_t$ the backward and $\hphi_t$ the forward heat flow of the two chosen functions.
In particular $\rho_0=\hphi\,\varphi_0$ and $\rho_1=\varphi\,\hphi_1$.
\end{keybox}
So a reference consists of six objects: $\varphi$, $\varphi_t$, $\hphi$, $\hphi_t$, $\rho_t$, $u$.
The two recipes differ only in \emph{which of these are chosen}, \emph{which are forced}, and
\emph{how the forced ones are computed}. In both, $\varphi$ is the Gaussian mixture
\eqref{eq:mixphi}, $\varphi=\sum_kw_k\Nc(\cdot;\mu_k,\Sigma_k)$, so that $\varphi_t$ and $u$
are closed form by the \textbf{Time potential} and \textbf{Drift} bullets of
Section~\ref{sec:gushchin}. Both scripts use the solver's convention $\sigma^2=2\eps$: wherever
$\sigma^2$ appears below the code writes \code{2*eps}, and the mixture parameters are the
arguments \code{(w, b, Q)} for $(w_k,\mu_k,\Sigma_k)$.

\paragraph{Recipe A: the paper route (\code{paper\_route\_sb.py}, class \code{FixedRho0SB}).}
This is Gushchin's construction, with the one object they never compute added at the end.
\begin{description}[leftmargin=1.5em,itemsep=2pt]
  \item[We fix two things.]
    \begin{enumerate}[itemsep=1pt]
      \item $\rho_0$: \emph{any} smooth positive density on $[0,1]^d$ that is negligible at the
            walls, normalized to mass one. It need not be a Gaussian or a mixture; the worked
            example in \code{candidate\_3d\_problems.tex} uses a ring,
            $\rho_0\propto\exp(-(|x-c|-r)^2/2w^2)$.
      \item $\varphi=\sum_kw_k\Nc(\cdot;\mu_k,\Sigma_k)$: the centres, widths and weights of
            the targets, chosen by hand (Section~\ref{sec:gushchin}). This is the paper's
            LSE potential $f^*=\sigma^2\log\varphi$.
    \end{enumerate}
  \item[Everything else is derived, in this order.]
    \begin{enumerate}[itemsep=1pt]
      \item $\varphi_t(x)=\sum_kw_k\Nc\big(x;\mu_k,\Sigma_k+\sigma^2(1-t)I\big)$, by the
            convolution rule. Closed form. In particular $\varphi_0$, which is also the paper's
            normalizer $Z_x$ up to a constant.
      \item $\hphi:=\rho_0/\varphi_0$. \emph{Forced}: this is the only choice that makes the
            $x$-marginal of $\hphi\Kc_\sigma\varphi$ equal to $\rho_0$ (first Schr\"odinger
            equation). Pointwise closed form, but not a mixture: for the ring it is a crescent,
            the side of the ring facing away from the targets, where $\varphi_0$ is smallest.
            The paper never writes this function down because its sampler never needs it.
      \item $\hphi_t=\Nc(0,\sigma^2tI)*\hphi$, the forward heat flow. \emph{The one numerical
            step.} Evaluate $\hphi$ on a fine $n_f^d$ grid; since the solver's cell centres form
            a tensor grid and the Gaussian kernel is separable, the convolution is
            $\hphi_t=F\times_1K_x\times_2K_y(\times_3K_z)$, with $F$ the fine-grid values and
            $K_x$ the $n\times n_f$ matrix of one-dimensional Gaussian weights times the fine
            cell width. Midpoint quadrature of a smooth Gaussian integrand converges
            spectrally: in the 2D example, $n_f=256$ is already below $10^{-9}$ relative to the
            peak at the earliest solver time, $t=1/32$. In 3D this is three tensor contractions,
            $O(n\,n_f^3)$ per time slice, seconds on a GPU.
      \item $\rho_t=\varphi_t\,\hphi_t$ and $\rho_1=\varphi\,\hphi_1$. No normalizing constant
            is needed: $\rho_0$ has mass one by construction and the Fokker--Planck equation
            conserves mass. $\rho_1$ is a product of a mixture and a numerically known
            function, so it is not explicit, but it is known at every cell to quadrature
            accuracy, which is all the solver needs.
      \item $u=\sigma^2\nabla\log\varphi_t$ from step~1, closed form and independent of
            $\rho_0$; $m=\rho_tu$.
      \item (Optional.) The static plan $\pi^*(y|x)=\varphi(y)\Nc(y;x,\sigma^2I)/\varphi_0(x)$,
            which is the paper's Proposition~3.3 mixture \eqref{eq:mixcond}.
    \end{enumerate}
  \item[What is exact and what is not.] Steps 1, 2, 5, 6 are closed form. Step~3 carries
    quadrature error, below $10^{-9}$ for $n_f\ge256$ in the example. Step~4 inherits it.
    The only other error is the $\R^d$-versus-box mismatch, bounded by the mass of $\rho_t$
    outside the walls ($8\cdot10^{-7}$ in the ring example). Both are orders of magnitude
    below the solver's $O(h^2)$ error.
  \item[Gain and loss.] Gain: complete freedom in $\rho_0$. Loss: $\rho_1$ has no formula, and
    there is no branch-mass table. Two practical rules: keep $\rho_0$ smooth on the scale
    $\sqrt{\sigma^2\,dt}$, the blur at the first time step, or $\hphi$ in step~2 is rough; and
    choose $\varphi$ so that $\varphi_0$ does not vary by orders of magnitude across the
    support of $\rho_0$, or $\hphi$ concentrates on a sliver of it (the ring-to-crescent
    effect, factor $261$ in the example).
\end{description}

\paragraph{Recipe B: both potentials as mixtures (\code{mixture\_sb.py}, class \code{MixtureSB}).}
The same construction with step~2 of Recipe~A replaced by a choice, which removes the numerical
step and makes everything explicit.
\begin{description}[leftmargin=1.5em,itemsep=2pt]
  \item[We fix two things, neither of which is a marginal.]
    \begin{enumerate}[itemsep=1pt]
      \item $\hphi=\sum_iu_i\Nc(\cdot;a_i,P_i)$: a Gaussian mixture for the \emph{initial}
            potential, with hand-picked $(u_i,a_i,P_i)$, the code's \code{(u, a, P)}. This
            replaces ``fix $\rho_0$''. $\rho_0$ is now an output.
      \item $\varphi=\sum_kw_k\Nc(\cdot;\mu_k,\Sigma_k)$ as in Recipe~A.
    \end{enumerate}
  \item[Everything else is derived.]
    \begin{enumerate}[itemsep=1pt]
      \item $\varphi_t(x)=\sum_kw_k\Nc\big(x;\mu_k,\Sigma_k+\sigma^2(1-t)I\big)$, as before.
      \item $\hphi_t(x)=\sum_iu_i\Nc\big(x;a_i,P_i+\sigma^2tI\big)$, by the convolution rule.
            This is Recipe~A's step~3, now closed form because $\hphi$ is a mixture.
      \item Mass. The product $\hphi\Kc_\sigma\varphi$ has total mass
            $K=\sum_{i,k}u_iw_k\,\Nc\big(a_i;\mu_k,P_i+\Sigma_k+\sigma^2I\big)$; divide by it.
            The normalized summands
            \[
              \Omega_{ik}=u_iw_k\,\Nc\big(a_i;\mu_k,P_i+\Sigma_k+\sigma^2I\big)\big/K,
              \qquad\textstyle\sum_{i,k}\Omega_{ik}=1,
            \]
            are the \emph{branch masses}: $\Omega_{ik}$ is the mass of the plan flowing from
            component $i$ of $\hphi$ to component $k$ of $\varphi$. They do not depend on $t$.
      \item $\rho_t=\varphi_t\hphi_t/K$. Each term is a product of two Gaussians in the same
            variable, so by \eqref{eq:gprod} with $A=P_i+\sigma^2tI$ and $B=\Sigma_k+\sigma^2(1-t)I$,
            \begin{equation}
              \label{eq:ourrho}
              \rho_t=\sum_{i,k}\Omega_{ik}\,\Nc\big(\cdot;\,c_{ik}(t),\,C_{ik}(t)\big),
            \end{equation}
            with $C_{ik}(t)=\big(A^{-1}+B^{-1}\big)^{-1}$ and
            $c_{ik}(t)=C_{ik}(t)\big(A^{-1}a_i+B^{-1}\mu_k\big)$.
            The constants $\Nc(a_i;\mu_k,A+B)$ produced by the product rule are exactly the
            numerators of $\Omega_{ik}$, since $A+B=P_i+\Sigma_k+\sigma^2I$ for every $t$. This is
            why the weights are $t$-independent. $\rho_0$ and $\rho_1$ are \eqref{eq:ourrho} at
            $t=0$ and $t=1$: both are mixtures of $I\!\cdot\!K$ Gaussians with the same weights.
      \item $u=\sigma^2\nabla\log\varphi_t$, $m=\rho_tu$, as before.
      \item (Optional.) The static plan $\pi^*(y|x)=\varphi(y)\Nc(y;x,\sigma^2I)/\varphi_0(x)$,
            unchanged, and the joint plan $\pi^*(x,y)=\sum_{i,k}\Omega_{ik}\times$ one
            bivariate Gaussian on $\R^{2d}$ per branch.
    \end{enumerate}
  \item[Why it all closes.] For one pair $(i,k)$,
    \[
      \hphi_i(x)\,\Kc_\sigma(x,y)\,\varphi_k(y)\;\propto\;
      \exp\!\Big(-\tfrac12\Big[(x-a_i)^{\!\top}P_i^{-1}(x-a_i)
      +\tfrac{|x-y|^2}{\sigma^2}+(y-\mu_k)^{\!\top}\Sigma_k^{-1}(y-\mu_k)\Big]\Big),
    \]
    whose exponent is one quadratic in the stacked vector $(x,y)\in\R^{2d}$. So each branch is a
    single Gaussian on $\R^{2d}$, the plan is a Gaussian mixture on $\R^{2d}$, and every marginal
    in $x$, in $y$, or at any intermediate time is a Gaussian mixture with the same weights.
    Example~\ref{ex:1d} is the $I=1$, $K=2$ case.
  \item[What is exact.] Everything; nothing is sampled or integrated. Checks recorded in the
    code and in \code{candidate\_3d\_problems.tex}: \eqref{eq:ourrho} equals the product
    $\varphi_t\hphi_t/K$ to $10^{-15}$; $\rho_1$ equals $\int\rho_0\pi^*(\cdot|x)\,dx$ by
    quadrature to $3\cdot10^{-13}$; the finite-difference Fokker--Planck residual decays as
    $O(h^2)$; and $I=K=1$ reproduces the Gaussian closed form of Bunne et al.
  \item[Gain and loss.] Gain: no numerical step, explicit $\rho_1$, and a branch-mass table
    $\Omega_{ik}$ that says in advance how much mass goes where. Loss: $\rho_0$ is a mixture
    of $I\!\cdot\!K$ Gaussians, not a free choice; a single-Gaussian $\rho_0$ with a
    multimodal $\rho_1$ is impossible here unless $K=1$. In the 1D example the two components of
    $\rho_0$ sit so close that it looks like one bump.
\end{description}

\paragraph{Side by side.}
\begin{center}\small
\begin{tabular}{@{}lll@{}}
\toprule
 & Recipe A (paper route) & Recipe B (both mixtures) \\
\midrule
chosen & $\rho_0$ (anything smooth), $\varphi$ (mixture) & $\hphi$ (mixture), $\varphi$ (mixture) \\
forced & $\hphi=\rho_0/\varphi_0$ & $\rho_0=\hphi\varphi_0/K$ \\
$\hphi_t$ & one Gaussian convolution, $\le10^{-9}$ & closed form \\
$\rho_1$ & numerical product, no formula & mixture, \eqref{eq:ourrho} at $t=1$ \\
$\rho_t$, $u$, $m$ & product and closed form & closed form \\
branch masses & not available & $\Omega_{ik}$, explicit \\
use when & a specific non-mixture $\rho_0$ is wanted & an all-explicit multimodal test is wanted \\
\bottomrule
\end{tabular}
\end{center}

\paragraph{What the solver is compared against.}
In the solver's convention, with $\eps=\sigma^2/2$:
\begin{itemize}
  \item $\rho_t$ on cell centres, from \eqref{eq:ourrho} (B) or from step~4 (A);
  \item $m=2\eps\,\hphi_t\nabla\varphi_t$ on the staggered faces (\code{staggered\_reference}
        in both classes evaluates at exactly the solver's points, so no interpolation);
  \item the cost $\int_0^1\!\!\int|m|^2/2\rho$, which by \eqref{eq:girsanov} equals
        $\sigma^2\KL(\mathbb{P}^*\|\mathbb{W}^\sigma)$; it is evaluated by quadrature of the
        reference on a fine grid, not by a formula of the Bunne type.
\end{itemize}

\paragraph{What this can and cannot do, and what is the paper's.}
The pair $(\rho_0,\rho_1)$ produced by Recipe~B is always two Gaussian mixtures, but a
\emph{derived} pair: shared weights $\Omega_{ik}$, means and covariances tied by
\eqref{eq:ourrho}. Given two arbitrary Gaussian mixtures, the bridge between them is not closed
form (the potentials are not mixtures), and neither recipe solves it. Both are generators of
test pairs with an exact solution, which is what a convergence study needs and what the paper
itself set out to build. From the paper we use three things: the idea of choosing the
potential and reading off the pair, which is what gives a multimodal bridge with an exact
answer at all (the single-Gaussian references of Bunne and Janati are the $I=K=1$ corner);
Theorem~3.2 as the citable optimality and uniqueness statement (the underlying fact is the
classical Schr\"odinger system); and the LSE formulas of Proposition~3.3 and Corollary~3.5,
which after the dictionary are \eqref{eq:mixcond} and the drift we evaluate. We do not use
the sampler, Theorem~3.1, or the benchmark pairs. Recipe~A adds the forward heat flow of
$\hphi$, which the paper never needs; Recipe~B adds the choice of $\hphi$ as a mixture.
Theorem~3.2 never required $P_0$ to be prescribed; the paper fixed it because they wanted a
given dataset or a clean Gaussian at $t=0$. Both recipes could be derived from the
Schr\"odinger system and the two Gaussian identities alone; citing the paper buys the
recognized framing and the theorem.

\paragraph{Boundary conditions.}
Every formula above is on $\R^d$. The solver lives on $[0,1]^d$ with no-flux walls, whose
reference process is \emph{reflected} Brownian motion and whose heat kernel is the Neumann
kernel, a cosine series rather than \eqref{eq:heat}. The references remain valid as long as
$\rho_t$ carries negligible mass at the walls for all $t$ (tail mass $\lesssim10^{-7}$ in the
examples), because then the two kernels agree to that accuracy on the support. The same
caveat applies to the Gaussian closed form of Bunne et al.

\section{Symbol table and reading map}
\label{sec:table}

\begin{center}
\small
\begin{tabular}{@{}p{3.9cm}p{3.1cm}p{3.0cm}p{3.6cm}@{}}
\toprule
Object & This primer & Gushchin et al.\ (check pdf) & Our code \\
\midrule
noise of reference process & $\sigma^2$ & $\epsilon$ & $2\eps$ \\
Fokker--Planck diffusion & $\sigma^2/2$ & $\epsilon/2$ & $\eps$ \\
EOT regularization, cost $\tfrac12|x-y|^2$ & $\sigma^2$ & $\epsilon$ & $2\eps$ \\
EOT regularization, cost $|x-y|^2$ (Janati) & $2\sigma^2$ & -- & $4\eps$ \\
source, target & $\rho_0,\rho_1$ & $P,Q$ & \code{rho0, rho1} \\
plan, conditionals & $\pi^*$, $\pi^*(\cdot|x)$ & $\pi^*$, $\pi^*(\cdot|x)$ & static plan \\
dual potential & $f=\epsilon\log\varphi$ (unused) & $f^*$ & -- \\
mixture component & $\Nc(\cdot;\mu_k,\Sigma_k)$, weight $w_k$ & $Q(\cdot|b_n,\epsilon^{-1}A_n)$, weight $w_n$ ($=c_k$ in the dictionary) & \code{(w,b,Q)}, \code{(u,a,P)} \\
conditional covariance, mean & $S_k$, $m_k(x)$ & $\Sigma_n=\epsilon(A_n+I)^{-1}$, $\mu_n(x)$ & -- \\
responsibilities & $\omega_k(t,x)$ & $\gamma_n$ (at $t=0$) & -- \\
potentials & $\hphi,\varphi$ & $\tilde\varphi^*,\ \varphi^*=e^{f^*/\epsilon}$ & \code{phihat, phi} \\
heat / entropic kernel & $p_1=\Nc(y;x,\sigma^2I)$ & $\exp(-|x-y|^2/2\epsilon)$ & -- \\
time-$t$ marginal & $\rho_t=\varphi_t\hphi_t$ & same & \code{rho(t)} on centres \\
drift, momentum & $u=\sigma^2\nabla\log\varphi_t$, $m=\rho u$ & $\epsilon\nabla\log\varphi_t$ & $m=2\eps\hphi_t\nabla\varphi_t$ on faces \\
dynamic cost & $\int\!\!\int|m|^2/2\rho=\sigma^2\KL(\mathbb{P}^*\|\mathbb{W}^\sigma)$ & $\epsilon\,\KL(\cdot\,\|\,W^\epsilon)$ & LADMM objective \\
\bottomrule
\end{tabular}
\end{center}

\paragraph{Reading map for Gushchin et al., Sections 2--3.}
\begin{enumerate}
  \item Section 2, entropic OT and weak OT: Sections~\ref{sec:couplings}--\ref{sec:weak} here.
        Goal: recognize \eqref{eq:eot}, \eqref{eq:plan} and the pointwise condition \eqref{eq:weakopt}.
  \item Section 2, dynamic Schr\"odinger bridge: Sections~\ref{sec:bm}--\ref{sec:sb}.
        Goal: $\rho_t=\varphi_t\hphi_t$ and drift $\epsilon\nabla\log\varphi_t$.
  \item Section 3.1, Theorem 3.2 and the drift display: Proposition~\ref{prop:construct}.
  \item Proposition 3.3 and Corollary 3.5: \eqref{eq:mixcond} and Example~\ref{ex:1d}.
  \item Stop there. Sections 4--6 are the benchmark and neural solvers.
\end{enumerate}

\paragraph{Further reading, in order of usefulness.}
L\'eonard, \emph{A survey of the Schr\"odinger problem and some of its connections with
optimal transport}, DCDS-A 34 (2014), arXiv:1308.0215, Sections 1--3, for the dynamic/static
reduction and the potentials; Peyr\'e \& Cuturi, \emph{Computational Optimal Transport},
Chapter 4, for entropic OT and Sinkhorn; Chen, Georgiou, Pavon, \emph{Stochastic control
liaisons}, SIAM Review 63 (2021), for the control form \eqref{eq:girsanov}; Bunne et al.,
AISTATS 2023, Theorem 3, and Janati et al., NeurIPS 2020, Theorem 1, for the Gaussian
closed forms used in Section~\ref{sec:ot}.

\end{document}
